\documentclass[11pt]{article}

\usepackage[utf8]{inputenc}
\usepackage[T1]{fontenc}
\usepackage{geometry}
\usepackage{amsmath}

\geometry{letterpaper, margin=2.5cm}

\title{Informe del módulo de recuperación}
\author{CSDI RAG Engine}
\date{}

\begin{document}

\maketitle

\section{Módulo de recuperación}

El módulo de recuperación es el componente central del sistema. Su responsabilidad es, dado el
texto libre de una consulta, identificar y ordenar los fragmentos del corpus que mayor relevancia
tienen para responderla. El módulo no es un recuperador único: es una cadena de etapas
especializadas que transforma la consulta original en un conjunto de fragmentos de alta calidad
listos para ser entregados al generador de respuestas.

Para que un sistema RAG funcione correctamente, no es suficiente recuperar documentos que
compartan palabras con la consulta: es necesario recuperar los fragmentos más relevantes
independientemente de cómo esté formulada la pregunta, del idioma en que esté escrita y de si
los términos exactos aparecen o no en el corpus. BM25 aplicado de forma aislada no puede
satisfacer esas condiciones simultáneamente. Un recuperador puramente léxico como BM25 exige
que los términos de la consulta coincidan con los del documento tras normalización morfológica;
falla ante paráfrasis, consultas en un idioma distinto al del corpus o preguntas conceptuales
sin vocabulario técnico exacto. Un recuperador puramente vectorial resuelve esas brechas
lingüísticas pero pierde precisión ante identificadores exactos como nombres de funciones,
mensajes de error o versiones de API, que los modelos de embeddings proyectan cerca de términos
relacionados aunque la distinción exacta sea la que el usuario necesita. Ninguno de los dos
enfoques por sí solo es suficiente para un sistema RAG de calidad en documentación técnica
multilingüe. El módulo de recuperación resuelve esta tensión combinando ambas señales y
añadiendo una fase de reordenamiento de alta precisión.

\paragraph{Diseño general del flujo de recuperación.}
El flujo completo de recuperación se organiza en tres etapas ordenadas: \emph{expansión opcional
de la consulta}, \emph{recuperación híbrida} y \emph{reordenamiento}. Esta organización no es
arbitraria: cada etapa actúa sobre la salida de la anterior y permite que los fragmentos
entregados al generador sean los más relevantes que el corpus puede ofrecer.

La expansión transforma la consulta original en una representación más rica antes de hacer la
búsqueda vectorial. La recuperación híbrida combina dos señales complementarias, léxica y
semántica, en una lista única de candidatos. El reordenamiento evalúa cada par consulta-fragmento
de forma conjunta para mejorar la precisión del ordenamiento final.

\paragraph{Componentes del recuperador léxico: BM25.}
El primer componente del recuperador híbrido es \texttt{BM25Retriever}. Este componente carga en
memoria, al arrancar la aplicación, el índice invertido construido por el módulo de indexación:
el diccionario de términos, las listas de postings, las longitudes documentales y las estadísticas
globales del corpus. Cuando recibe una consulta, aplica exactamente la misma tokenización usada
durante la indexación, consulta las listas de postings de los términos presentes y calcula la
puntuación BM25 de cada documento candidato.

El recuperador léxico aporta una ventaja concreta en el dominio del sistema. La documentación
técnica contiene nombres de funciones, clases, mensajes de error, operadores y palabras clave
del lenguaje que rara vez son parafraseados: quien busca \texttt{asyncio.gather} o
\texttt{Array.prototype.flatMap} espera encontrar exactamente esas cadenas en los documentos
recuperados. BM25 aprovecha esa propiedad porque compara términos directamente después de
normalización morfológica, sin proyectarlos a un espacio de representación que podría absorber
distinciones relevantes.

\paragraph{Componentes del recuperador semántico: embeddings y FAISS.}
El segundo componente es \texttt{VectorRetriever}. Utiliza el modelo de embeddings
\texttt{BAAI/bge-m3}, un modelo multilingüe que produce vectores de 1024 dimensiones. Esta
elección no fue la inicial: durante el desarrollo se trabajó primero con un modelo de un solo
idioma que producía vectores de 384 dimensiones. Las pruebas con ese modelo mostraron que la
calidad de la recuperación caía de forma notable cuando el usuario formulaba consultas en
español sobre documentación escrita en inglés, ya que el modelo no tenía representaciones
compartidas entre ambos idiomas. Al cambiar a \texttt{BAAI/bge-m3}, que fue entrenado sobre
datos en más de cien idiomas, esa brecha desapareció: la consulta en español y el fragmento
en inglés quedan proyectados a regiones cercanas del espacio de representación cuando
describen el mismo concepto. El costo de esta mejora es un vector tres veces más grande
(1024 frente a 384 dimensiones), lo que aumenta el uso de memoria del índice y el tiempo de
codificación, pero se consideró un intercambio aceptable dado el perfil multilingüe real del
sistema.

Los vectores de los fragmentos se almacenan en un índice FAISS de tipo HNSW
(\emph{Hierarchical Navigable Small World}). HNSW es una estructura de grafo jerárquico que
permite realizar búsqueda aproximada de vecinos más cercanos con muy bajo costo de consulta.
A diferencia de índices planos que comparan el vector de consulta contra todos los vectores del
corpus, HNSW navega capas del grafo reduciendo el número de comparaciones necesarias. Los
parámetros de construcción \texttt{HNSW\_M=32} y \texttt{HNSW\_EF\_CONSTRUCTION=200}
controlan la densidad del grafo y la precisión de la indexación; el parámetro de búsqueda
\texttt{HNSW\_EF\_SEARCH=50} controla el equilibrio entre velocidad y exhaustividad en tiempo
de consulta.

La métrica de similaridad utilizada es el producto interno, que equivale a similaridad coseno
cuando los vectores están normalizados en norma L2. Los embeddings se normalizan al momento de
la indexación y al momento de codificar la consulta, garantizando coherencia entre ambas fases.
Cuando se recibe una consulta, \texttt{VectorRetriever} la codifica, consulta el índice FAISS
y filtra automáticamente los documentos marcados como eliminados antes de devolver resultados.

\paragraph{Elección de HNSW sobre índices alternativos.}
La decisión de usar HNSW en lugar de un índice plano o una estructura basada en árboles
responde a la escala esperada del corpus. Un índice plano en FAISS compara el vector de consulta
contra todos los vectores en tiempo lineal, lo cual es exacto pero impracticable cuando el corpus
tiene decenas o cientos de miles de fragmentos. Los índices de cuantización como IVF reducen ese
costo dividiendo el espacio en celdas, pero requieren entrenamiento previo sobre el corpus y
pueden perder calidad en corpus pequeños o poco densos.

HNSW no requiere entrenamiento, escala de forma eficiente con el tamaño del corpus y mantiene
una alta tasa de recuperación en la práctica. Para el dominio de documentación técnica, donde la
precisión en los primeros resultados es crítica, esta estructura es más adecuada que alternativas
que sacrifican calidad para ganar velocidad.

\paragraph{Fusión híbrida mediante Reciprocal Rank Fusion.}
El componente central del módulo es \texttt{HybridRetriever}. Este componente ejecuta
\texttt{BM25Retriever} y \texttt{VectorRetriever} de forma concurrente usando un
\texttt{ThreadPoolExecutor}, reúne ambas listas de resultados y las combina mediante
\emph{Reciprocal Rank Fusion} (RRF).

RRF no utiliza directamente las puntuaciones de cada recuperador, sino únicamente el rango de
cada documento dentro de su lista. La puntuación combinada de un documento $d$ se calcula como:

\[
\mathrm{score}_{\mathrm{RRF}}(d) =
\sum_{r \in R}
w_r \cdot \frac{1}{k + \mathrm{rank}_r(d)}
\]

donde $R$ es el conjunto de recuperadores activos, $w_r$ es el peso asignado al recuperador $r$,
$\mathrm{rank}_r(d)$ es la posición de $d$ en la lista del recuperador $r$ (empezando en 1), y
$k$ es una constante de suavizado con valor por defecto $k=60$. Si un documento no aparece en
la lista de un recuperador, no contribuye a la suma por ese recuperador.

Los pesos por defecto son \texttt{bm25\_weight=0.3} y \texttt{vector\_weight=0.7}. Esta
distribución reconoce que la similitud semántica es la señal más general, pero no descarta el
aporte de la coincidencia léxica para términos técnicos. Ambos pesos son configurables en tiempo
de ejecución mediante \texttt{update\_weights()}, lo que permite ajustar el comportamiento del
recuperador sin reiniciar la aplicación.

\paragraph{Por qué RRF y no combinación directa de puntuaciones.}
Combinar directamente las puntuaciones de BM25 y del recuperador vectorial plantea un problema
de escala. BM25 produce puntuaciones cuya magnitud depende del tamaño del corpus, la frecuencia
de los términos y la longitud de los documentos; los scores de FAISS producen valores de producto
interno que dependen de la distribución del modelo de embeddings. Estas magnitudes son
heterogéneas y no comparables sin una normalización previa que puede ser sensible a la
distribución específica de cada consulta.

RRF evita este problema por completo: solo considera el orden, no el valor. Un documento que
aparece bien posicionado en ambas listas obtendrá una puntuación alta independientemente de la
escala de cada recuperador. Además, RRF es robusto ante casos donde un recuperador devuelve
pocos resultados o donde las puntuaciones están muy concentradas en pocos documentos. Por estas
razones, RRF es una elección técnicamente sólida para la fusión híbrida sin requerir calibración
por corpus.

\paragraph{Expansión de consulta mediante HyDE.}
Para mejorar la calidad de la búsqueda vectorial en consultas breves o abstractas, el sistema
implementa \emph{Hypothetical Document Embeddings} (HyDE). Cuando está habilitado, antes de
codificar la consulta con el modelo de embeddings, el pipeline solicita al modelo generativo que
produzca un documento hipotético que respondería la consulta. Ese documento hipotético, en lugar
del texto de la consulta original, se usa como entrada al modelo de embeddings para la búsqueda
vectorial.

La motivación técnica es que las consultas y los documentos tienen distribuciones textuales
distintas: una consulta es corta y directiva mientras que un documento técnico es extenso y
descriptivo. El modelo de embeddings produce vectores para ambos, pero un documento hipotético
está más cerca, en ese espacio de representación, a los documentos reales del corpus que la
consulta original. HyDE se aplica únicamente sobre el componente vectorial; BM25 continúa
operando sobre la consulta original porque la coincidencia léxica no se beneficia de esta
expansión.

Si la expansión falla por cualquier razón, el sistema retrocede a la consulta original sin
interrumpir el flujo.

HyDE está desactivado por defecto por una razón de diseño que va más allá de la latencia: el
sistema establece como principio que el modelo generativo interviene únicamente en la etapa de
generación de respuesta. Usar el LLM dentro del pipeline de recuperación rompe esa separación,
convirtiendo una etapa que debería ser determinista y auditable en una etapa dependiente de un
modelo externo. Esto introduce un punto adicional de fallo, aumenta el costo por consulta y
hace más difícil aislar la causa cuando la recuperación produce resultados incorrectos.

Sin embargo, HyDE se mantiene en el sistema porque ilustra de forma concreta cómo una
intervención mínima del LLM —generar un único párrafo antes de la búsqueda vectorial— puede
mejorar la calidad de la recuperación sin modificar ningún componente del índice ni del
pipeline. Es un ejemplo observable de la tensión entre pureza arquitectónica y rendimiento
práctico, y su presencia como opción configurable permite evaluar esa diferencia en el corpus
real. Se activa mediante la variable de entorno \texttt{HYDE\_ENABLED=true}.

\paragraph{Reordenamiento mediante Cross-Encoder.}
La recuperación híbrida produce una lista de candidatos ordenados por su puntuación RRF. Sin
embargo, el ranking RRF se basa en el orden producido por dos modelos de primera etapa que
procesan consulta y documento de forma independiente. Un \emph{cross-encoder} puede evaluar cada
par (consulta, fragmento) conjuntamente, capturando interacciones entre ellos que los encoders
independientes no modelan.

El reordenador implementado es \texttt{CrossEncoderReranker}, que utiliza el modelo
\texttt{cross-encoder/ms-marco-MiniLM-L-6-v2}. La elección de este modelo responde a tres
criterios. Primero, fue entrenado sobre MS-MARCO, un benchmark de 8,8 millones de pares
(consulta, pasaje) extraídos de búsquedas web reales, que es el conjunto de entrenamiento
estándar para sistemas de ranking de pasajes. Segundo, la arquitectura MiniLM de 6 capas
ofrece un compromiso entre calidad y latencia: con aproximadamente 22 millones de parámetros,
puede puntuar los 30 pares en CPU sin introducir una demora apreciable. Tercero, cross-encoders
más grandes, como \texttt{ms-marco-electra-base}, ofrecen mayor precisión de ranking pero
multiplican el tiempo de inferencia, lo que no está justificado para un conjunto de 30
candidatos ya preseleccionados por la recuperación híbrida.

En la primera etapa del pipeline se recuperan \texttt{reranker\_candidate\_k=30} candidatos;
el reordenador puntúa los 30 pares (consulta, fragmento) y devuelve los
\texttt{context\_chunks=15} mejor posicionados. La relación 2:1 entre candidatos y fragmentos
finales es una proporción deliberada: recuperar menos candidatos reduciría la cobertura del
corpus expuesta al cross-encoder; recuperar más incrementaría el costo de inferencia sin
garantizar una mejora proporcional en la calidad del conjunto final.

El uso de una primera etapa de candidatos amplia seguida de un reordenador preciso es una
estrategia estándar en recuperación de información moderna. El cross-encoder es computacionalmente
más costoso que los encoders de primera etapa, por lo que no puede aplicarse sobre todo el corpus
directamente, pero puede aplicarse eficientemente sobre un conjunto reducido de candidatos ya
filtrados.

\paragraph{Por qué BM25 solo no es suficiente para un sistema RAG.}
Un sistema RAG exige que el módulo de recuperación entregue al generador los fragmentos con
mayor densidad de información relevante para la consulta. Si esos fragmentos no son los
correctos, el generador producirá una respuesta incorrecta o incompleta sin importar su
capacidad. BM25 por sí solo no puede garantizar esa calidad de recuperación en las condiciones
reales del sistema por tres razones estructurales.

Primero, BM25 es sensible al vocabulario: una consulta formulada en español sobre documentación
en inglés no encuentra coincidencias léxicas en el corpus, devolviendo resultados vacíos o
irrelevantes aunque los documentos necesarios existan. Segundo, BM25 no captura relaciones
semánticas: una consulta como \emph{``cómo iterar sobre una lista''} no recuperará un fragmento
que explica \texttt{forEach} si ese término no aparece en la consulta. Tercero, incluso cuando
BM25 recupera los documentos correctos, los ordena por frecuencia de término, no por relevancia
semántica al contexto completo de la consulta, lo que degrada la calidad del conjunto enviado
al generador.

La recuperación vectorial resuelve el primer y segundo problema pero introduce el suyo propio:
colapsa la distinción entre identificadores exactos. La fusión híbrida resuelve ambos problemas
a la vez. El reordenamiento con cross-encoder corrige el tercero. Cada componente es necesario
para cumplir el requisito de un sistema RAG funcional; omitir cualquiera de ellos introduce
una categoría de fallo que los demás no pueden compensar.

\paragraph{Integración en el pipeline RAG.}
El módulo de recuperación está integrado en \texttt{RAGPipeline}, que coordina la secuencia
completa desde la consulta del usuario hasta la respuesta generada. El pipeline invoca primero
a \texttt{HybridRetriever} para obtener candidatos, opcionalmente a través de la expansión HyDE.
Si el reordenador está activo, reordena los candidatos antes de seleccionar los fragmentos
finales. Con esos fragmentos construye el prompt y lo envía al modelo generativo.

El resultado devuelto al cliente incluye la respuesta generada, las fuentes utilizadas con sus
URLs y títulos, el modelo utilizado y los conteos de tokens. Esta trazabilidad es parte del
diseño del módulo: el usuario puede verificar de dónde proviene la información y el sistema
puede auditar qué fragmentos influyeron en cada respuesta.

\paragraph{API del módulo de recuperación.}
El módulo expone varios endpoints dentro del servidor. El endpoint
\texttt{POST /api/v1/search} ejecuta la búsqueda híbrida directamente y devuelve los fragmentos
recuperados con sus metadatos de fuente, título, breadcrumb y texto. El endpoint
\texttt{POST /api/v1/search/bm25} expone el recuperador léxico de forma aislada, útil para
diagnóstico y comparación. El endpoint \texttt{POST /api/v1/vector/search} expone el recuperador
vectorial de forma aislada. El endpoint \texttt{POST /api/v1/rag/query} ejecuta el pipeline
completo incluyendo generación de respuesta.

La existencia de endpoints individuales para cada componente del recuperador permite validar de
forma independiente cada señal de relevancia y diagnosticar casos donde un enfoque aporta más
que el otro. Esto fue una decisión de diseño tomada durante el desarrollo para facilitar la
depuración y el ajuste de pesos de fusión.

\paragraph{Configuración y parámetros clave.}
Los parámetros del módulo de recuperación se controlan mediante variables de entorno y son
ajustables en tiempo de ejecución a través de la API de configuración. Los parámetros más
relevantes son el número de candidatos de primera etapa para el reordenador
(\texttt{RERANKER\_CANDIDATE\_K=30}), el número de fragmentos finales enviados al generador
(\texttt{LLM\_CONTEXT\_CHUNKS=15}), los pesos de fusión BM25 y vectorial
(\texttt{bm25\_weight=0.3}, \texttt{vector\_weight=0.7}) y el parámetro de suavizado de RRF
($k=60$).

Este diseño de configuración externalizada permite adaptar el comportamiento del recuperador a
condiciones cambiantes del corpus sin modificar el código. Si el corpus crece o cambia de
dominio, los pesos de fusión y el tamaño del conjunto de candidatos pueden ajustarse de forma
controlada y sin interrumpir el servicio.

\paragraph{Justificación de los valores de los parámetros.}
Los valores por defecto de los parámetros del módulo no fueron elegidos arbitrariamente. Cada
uno responde a una propiedad del dominio o a evidencia empírica de la literatura.

El parámetro de suavizado de RRF $k=60$ proviene directamente del trabajo original de Cormack
et al.\ (2009), que demostró que valores en el rango 40--80 son estables ante variaciones del
corpus; valores más bajos amplifican diferencias de rango en la cabeza de la lista, y valores
más altos las suavizan en exceso. Los pesos de fusión $w_{\text{BM25}}=0{,}3$ y
$w_{\text{vector}}=0{,}7$ reflejan que el corpus incluye documentación en inglés consultada
con frecuencia en español: el componente vectorial, respaldado por el modelo multilingüe
\texttt{BAAI/bge-m3}, tiene mayor capacidad de resolver esa brecha lingüística y por eso
recibe mayor peso. El componente léxico retiene peso suficiente para no perder la ventaja de
BM25 ante identificadores exactos.

Los parámetros de BM25 $k_1=1{,}5$ y $b=0{,}75$ corresponden a los valores canónicos
recomendados por Robertson \& Zaragoza (2009) para colecciones de documentación técnica, donde
la frecuencia de término importa pero los documentos tienen longitudes heterogéneas. El tamaño
de fragmento de 256 palabras con 32 de solapamiento fue elegido porque mantiene suficiente
contexto para que el generador entienda el pasaje, sin superar el rango en el que los modelos
de embeddings empiezan a perder coherencia semántica al codificar fragmentos largos.

Los parámetros de HNSW ($M=32$, $\mathit{ef\_construction}=200$, $\mathit{ef\_search}=50$)
corresponden a la configuración de alta calidad del índice: $M=32$ produce un grafo denso que
mantiene alta tasa de recuperación, y $\mathit{ef\_construction}=200$ asegura que la
construcción del grafo explore suficientes vecinos para no crear atajos subóptimos. El valor
$\mathit{ef\_search}=50$ equilibra exhaustividad y velocidad en tiempo de consulta.

\paragraph{Justificación técnica de la solución.}
El módulo de recuperación está diseñado para superar las deficiencias concretas de los
recuperadores simples. Identificar esas deficiencias es necesario para entender por qué cada
componente del módulo existe.

Un recuperador puramente léxico como BM25 falla por \emph{desajuste de vocabulario}: si el
usuario formula la consulta en español y la documentación está en inglés, o si usa sinónimos o
reformulaciones del término técnico presente en el documento, BM25 no encuentra ninguna
coincidencia y devuelve resultados vacíos o irrelevantes. Un recuperador puramente vectorial
falla por \emph{colapso de distinción exacta}: los modelos de embeddings proyectan términos
semánticamente próximos a regiones cercanas del espacio, lo que hace que \texttt{Array.flat()}
y \texttt{Array.flatMap()} aparezcan como casi equivalentes aunque el usuario esté buscando
exactamente uno de los dos. Combinar ambos enfoques mediante RRF resuelve estos dos problemas
sin requerir normalización de puntuaciones, porque opera únicamente sobre rangos.

Incluso con recuperación híbrida, el problema del \emph{ranking por encoders independientes}
persiste. Tanto BM25 como el modelo de embeddings puntúan consulta y fragmento por separado, sin
modelar la interacción entre ellos. Un fragmento que contiene todos los términos clave puede
estar mal posicionado si esos términos no aparecen en el contexto semántico correcto. El
cross-encoder resuelve esto evaluando cada par (consulta, fragmento) de forma conjunta; su costo
computacional se controla aplicándolo únicamente sobre los 30 candidatos preseleccionados por la
recuperación híbrida.

La expansión HyDE aborda un cuarto problema: la \emph{asimetría distribucional entre consultas
y documentos}. Una consulta es corta y directiva; un documento técnico es extenso y descriptivo.
El modelo de embeddings produce un vector para cada uno, pero esos vectores parten de
distribuciones textuales distintas y no siempre quedan cerca en el espacio de representación
aunque describan el mismo concepto. HyDE reduce esta brecha generando, antes de la búsqueda
vectorial, un párrafo hipotético con la forma de un documento real que respondería la consulta.
Ese párrafo es lo que se codifica y se usa como vector de consulta; BM25 continúa operando
sobre la consulta original porque la coincidencia léxica no se beneficia de esta transformación.

HyDE es un mecanismo de refuerzo opcional dentro del sistema. Está desactivado por defecto
porque introduce una llamada adicional al modelo generativo antes de cada recuperación,
incrementando la latencia de la consulta. Se activa mediante la variable de entorno
\texttt{HYDE\_ENABLED=true} en escenarios donde la calidad de la recuperación vectorial es
prioritaria sobre la velocidad. Si la llamada al modelo falla por cualquier razón, el sistema
retrocede a la consulta original sin interrumpir el flujo ni degradar los otros componentes.
HyDE no reemplaza la recuperación híbrida ni el reordenamiento: los complementa en la etapa de
codificación vectorial, y su contribución es independiente de los demás.

La separación entre recuperación léxica, semántica, fusión y reordenamiento también mejora la
mantenibilidad del sistema. Cada componente tiene límites claros, puede probarse de forma
independiente y puede evolucionar sin afectar a los demás. Esta arquitectura ofrece una base
sólida para la integración de futuras mejoras, como modelos de embeddings más capaces,
estrategias alternativas de fusión o mecanismos de retroalimentación sobre la calidad de los
resultados.

Las decisiones de diseño descritas en este módulo están respaldadas por evidencia publicada.
Cormack et al.\ (2009) demostraron que RRF supera la combinación directa de puntuaciones y los
métodos de aprendizaje de ranking en una variedad de colecciones sin requerir calibración.
Malkov \& Yashunin (2020) establecieron HNSW como el estándar para búsqueda aproximada de
vecinos en alta dimensión, superando en velocidad y tasa de recuperación a los índices planos
y a las variantes IVF de FAISS para corpus de tamaño medio. Gao et al.\ (2023) introdujeron
HyDE y mostraron que mejora la recuperación densa en escenarios de \emph{zero-shot} sin
requerir datos de relevancia etiquetados, que es exactamente la condición del sistema. El
modelo \texttt{BAAI/bge-m3} ocupa consistentemente los primeros lugares en el benchmark MTEB
(\emph{Massive Text Embedding Benchmark}) para tareas de recuperación multilingüe.

\paragraph{Costo computacional por etapa.}
El pipeline de recuperación acumula costo en cada etapa. La siguiente tabla resume el costo
relativo de cada componente y si es posible desactivarlo:

\begin{center}
\begin{tabular}{llll}
\hline
\textbf{Etapa} & \textbf{Complejidad} & \textbf{Costo relativo} & \textbf{Opcional} \\
\hline
BM25 & $O(q \cdot t)$ en memoria & Bajo & No \\
Vectorización HNSW & $O(\log n)$ aprox. & Bajo & No \\
RRF (fusión) & $O(c)$ sobre candidatos & Despreciable & No \\
Cross-encoder (30 pares) & $O(c \cdot L)$ por par & Medio & Sí (ENV) \\
HyDE (llamada al LLM) & Una inferencia generativa & Alto & Sí (off por defecto) \\
\hline
\end{tabular}
\end{center}

En la configuración por defecto, las etapas costosas en orden son el cross-encoder (activo) y
la vectorización de la consulta. HyDE, al estar desactivado por defecto, no contribuye al
costo nominal del pipeline. El cross-encoder opera sobre un máximo fijo de 30 pares
independientemente del tamaño del corpus, lo que mantiene su costo acotado y predecible.

\paragraph{Limitaciones del módulo.}
El diseño del módulo tiene limitaciones conocidas que acotan su alcance.

El tokenizador de BM25 aplica \emph{stopwords} y \emph{stemming} en inglés
(\texttt{SnowballStemmer}, \texttt{language='english'}). Consultas formuladas exclusivamente
en español son procesadas correctamente por el componente vectorial, que utiliza un modelo
multilingüe, pero el componente léxico pierde eficiencia porque el stemmer no normaliza raíces
en español y las stopwords en español no son filtradas. En la práctica, el peso vectorial
dominante ($w=0{,}7$) compensa esta limitación, pero documentos con términos técnicos en
español podrían no ser recuperados por BM25 con la misma precisión que en inglés.

El modelo de cross-encoder \texttt{ms-marco-MiniLM-L-6-v2} fue entrenado exclusivamente sobre
datos en inglés. Al reordenar pares (consulta en español, fragmento en inglés) puede degradar
la posición de fragmentos relevantes que en inglés habrían sido bien posicionados. Este efecto
es difícil de cuantificar sin un conjunto de evaluación etiquetado en el dominio.

La expansión HyDE genera el documento hipotético sin ningún contexto del corpus indexado. Si el
modelo generativo produce un párrafo que describe un concepto inexistente en los documentos, el
vector resultante puede alejar la búsqueda de los fragmentos correctos. Este riesgo aumenta
cuanto más especializado es el corpus y menos representado está en los datos de entrenamiento
del LLM.

El tamaño de fragmento de 256 palabras está definido de forma estática en el código y no es
configurable mediante variables de entorno. Corpora con documentos de estructura muy diferente,
como páginas de referencia de API cortas o guías de instalación extensas, podrían beneficiarse
de tamaños de fragmento distintos.

\end{document}
